\documentclass[12pt]{article}
\usepackage{amsmath,amssymb,amsthm}
\usepackage[margin=1in]{geometry}
\usepackage{enumitem}
\usepackage{booktabs}
\usepackage{hyperref}

\newtheorem{theorem}{Theorem}[section]
\newtheorem{lemma}[theorem]{Lemma}
\newtheorem{corollary}[theorem]{Corollary}
\newtheorem{proposition}[theorem]{Proposition}
\newtheorem{conjecture}[theorem]{Conjecture}
\theoremstyle{definition}
\newtheorem{definition}[theorem]{Definition}
\theoremstyle{remark}
\newtheorem{remark}[theorem]{Remark}

\newcommand{\CC}{\mathbb{C}}
\newcommand{\RR}{\mathbb{R}}
\newcommand{\HH}{\mathbb{H}}
\newcommand{\OO}{\mathbb{O}}
\newcommand{\agut}{\alpha_{\mathrm{GUT}}}
\newcommand{\Tr}{\mathrm{Tr}}
\newcommand{\rank}{\mathrm{rank}}

\title{Born-Rule Weighted Gram Graphs,\\
  the Doubly Irreducible Number,\\
  and the Critical Line}
\author{Mingu Jeong\\Independent Researcher
\and Claude\\Anthropic}
\date{\today}

\begin{document}
\maketitle

\begin{abstract}
We prove that $\CC$ is the unique normed division algebra
where the statistical convergence boundary ($\sigma_{\mathrm{stat}}$)
and the geometric phase equipartition ($\sigma_{\mathrm{geom}}$)
of random Dirichlet series coincide, both equalling~$1/2$.
This coincidence is traced to the \emph{doubly irreducible}
nature of the number~2: it is simultaneously the smallest
additive atom and the unique extension degree over~$\RR$.
In the DRLT framework, $1/2 = 1/n_T = 1/c$ (inverse temporal
dimension = inverse lattice speed of light).

We then show that Born-rule weighted Gram graphs
($W_{ij} = |\langle\psi_i|\psi_j\rangle|^2$, no thresholding)
are Ramanujan for $N \leq N_c(d)$, with all Ihara zeta zeros on
$\mathrm{Re}(s) = 1/2$.
The Ramanujan ratio scales as $\sim A \cdot d^{-\alpha}$
($\alpha \approx 0.67$, $R^2 = 0.996$), giving a critical
dimension $d_c \approx 3$.
Since the DRLT chiral decomposition forces $d = 5 > d_c$,
the Gram graph lies in the Ramanujan-safe region.

No free parameters appear in any result.
\end{abstract}

%=====================================================================
\section{Introduction}
\label{sec:intro}
%=====================================================================

The Riemann Hypothesis (RH) asserts that all nontrivial zeros of
$\zeta(s)$ have $\mathrm{Re}(s) = 1/2$.
Two long-standing mysteries surround this assertion:
\begin{enumerate}[label=(\roman*)]
\item \emph{Why $1/2$?} What structural property of $\zeta(s)$
  fixes the critical line at this specific value?
\item \emph{Why GUE?} Montgomery--Odlyzko showed that $\zeta$ zero
  statistics match the Gaussian Unitary Ensemble (GUE), but no
  framework explains \emph{why} GUE and not GOE or GSE.
\end{enumerate}

This paper answers both questions within the Dynamic Resolution
Lattice Theory (DRLT), building on Papers~1--4 in this
series~\cite{Paper1,Paper2,Paper3,Paper4}.
The answers are:
\begin{itemize}
\item[(i)] $1/2$ is the CLT convergence boundary for
  unit-modulus Dirichlet series, equal to $1/\dim_\RR(\CC)$
  because $\CC = \RR^2$.
\item[(ii)] GUE arises because $\CC$ forces Dyson index $\beta = 2$.
\end{itemize}
Crucially, (i)~is \emph{universal} (holds for $\RR$ too),
while (ii)~is \emph{$\CC$-specific}.
Their unification rests on the fact that $\dim_\RR(\CC) = 2 = n_T$,
the temporal sector dimension of the chiral decomposition.

We then construct Born-rule weighted Gram graphs and show they
are naturally Ramanujan for the DRLT dimension $d = 5$,
establishing a \emph{discrete Riemann Hypothesis} that holds
exactly for finite $N$ and softens at rate $O(1/\sqrt{N})$.

%=====================================================================
\section{The Doubly Irreducible Number}
\label{sec:doubly}
%=====================================================================

\begin{definition}
An integer $n \geq 2$ is an \emph{additive atom} if it admits
no partition $n = a + b$ with $a, b \geq 2$.
\end{definition}

\begin{definition}
An integer $n \geq 2$ is an \emph{extension atom over $\RR$} if
there exists an irreducible polynomial of degree $n$ in $\RR[x]$.
\end{definition}

\begin{lemma}\label{lem:add}
The set of additive atoms is $\{2, 3\}$.
\end{lemma}
\begin{proof}
For $n \geq 4$: $n = 2 + (n-2)$ with $n-2 \geq 2$.
For $n = 2$: every partition has a part $< 2$.
For $n = 3$: every partition has a part $< 2$.
\end{proof}

\begin{lemma}\label{lem:ext}
The set of extension atoms over $\RR$ is $\{2\}$.
\end{lemma}
\begin{proof}
By the Fundamental Theorem of Algebra, $\CC$ is algebraically
closed. Since $[\CC:\RR] = 2$, the only nontrivial irreducible
real polynomial degree is $2$ (e.g., $x^2 + 1$).
Every odd-degree polynomial has a real root (IVT), hence is
reducible.
\end{proof}

\begin{theorem}[Unique Doubly Irreducible Number]
\label{thm:doubly}
$$\{\text{additive atoms}\} \cap
  \{\text{extension atoms over } \RR\} = \{2,3\} \cap \{2\} = \{2\}$$
\end{theorem}

\begin{corollary}\label{cor:nT}
$n_T = \dim_\RR(\CC) = 2$.
Both equalities hold because $2$ is the unique doubly irreducible
number: $n_T = 2$ as the smallest additive atom (Paper~1), and
$\dim_\RR(\CC) = 2$ as the unique extension atom.
\end{corollary}

%=====================================================================
\section{The Two Boundaries Theorem}
\label{sec:boundaries}
%=====================================================================

\begin{definition}
For a normed division algebra $K$ over $\RR$ with $n_K = \dim_\RR(K)$:
\begin{enumerate}[label=(\alph*)]
\item The \emph{statistical boundary}
  $\sigma_{\mathrm{stat}}(K) :=
  \inf\{\sigma > 0 : \sum a_k/k^\sigma \text{ converges a.s.}\}$
  where $\{a_k\}$ are iid uniform on $U(K)$.
\item The \emph{geometric boundary}
  $\sigma_{\mathrm{geom}}(K) := \mathbb{E}[x_1^2]$
  for $(x_1, \ldots, x_{n_K})$ uniform on $S^{n_K - 1}$.
\end{enumerate}
\end{definition}

\begin{lemma}\label{lem:stat}
$\sigma_{\mathrm{stat}}(K) = 1/2$ for all $K \in \{\RR, \CC, \HH, \OO\}$.
\end{lemma}
\begin{proof}
$\mathrm{Var}(a_k/k^\sigma) = |a_k|^2/k^{2\sigma} = 1/k^{2\sigma}$.
By the Kolmogorov three-series theorem,
$\sum 1/k^{2\sigma} < \infty$ iff $2\sigma > 1$.
\end{proof}

\begin{lemma}\label{lem:geom}
$\sigma_{\mathrm{geom}}(K) = 1/n_K$.
\end{lemma}
\begin{proof}
By spherical symmetry,
$\mathbb{E}[x_i^2] = \mathbb{E}[x_j^2]$ for all $i, j$.
Since $\sum x_i^2 = 1$, we get $n_K \cdot \mathbb{E}[x_1^2] = 1$.
\end{proof}

\begin{theorem}[Two Boundaries]\label{thm:two}
$\sigma_{\mathrm{stat}}(K) = \sigma_{\mathrm{geom}}(K)$
if and only if $K = \CC$.
\end{theorem}
\begin{proof}
$1/2 = 1/n_K$ iff $n_K = 2$ iff $K = \CC$.

\begin{center}
\begin{tabular}{lcccc}
\toprule
$K$ & $n_K$ & $\sigma_{\mathrm{stat}}$ & $\sigma_{\mathrm{geom}}$ & Equal? \\
\midrule
$\RR$ & 1 & 1/2 & 1 & No \\
$\CC$ & \textbf{2} & \textbf{1/2} & \textbf{1/2} & \textbf{Yes} \\
$\HH$ & 4 & 1/2 & 1/4 & No \\
$\OO$ & 8 & 1/2 & 1/8 & No \\
\bottomrule
\end{tabular}
\end{center}
\vspace{-1em}
\end{proof}

\begin{corollary}[The 1/2 Chain]\label{cor:chain}
$$\frac{1}{2} = \frac{1}{n_T} = \frac{1}{\dim_\RR(\CC)}
  = \frac{1}{c} = \sigma_{\mathrm{stat}} = \sigma_{\mathrm{geom}}$$
\end{corollary}

%=====================================================================
\section{GUE from $\CC$}
\label{sec:gue}
%=====================================================================

The Dyson classification assigns index $\beta$ to random
matrix ensembles based on the underlying field~\cite{Dyson}:
$\beta = 1$ (GOE, $\RR$), $\beta = 2$ (GUE, $\CC$),
$\beta = 4$ (GSE, $\HH$).

\begin{proposition}\label{prop:gue}
The Gram matrix $G_{ij} = \langle\psi_i|\psi_j\rangle$
($\psi_i \in \CC^5$) belongs to the $\beta = 2$ universality
class.
\end{proposition}

\begin{proof}
$G = \Psi\Psi^\dagger$ is Hermitian over $\CC$.
By the Dyson classification, Hermitian matrices over $\CC$
have $\beta = 2$.
\end{proof}

Numerical confirmation (EXP\_071, ratio statistic~\cite{Atas2013}):
\begin{itemize}
\item $\CC^5$: $\langle r \rangle = 0.594 \pm 0.002$
  (GUE theory: 0.603)
\item $\RR^5$: $\langle r \rangle = 0.526 \pm 0.002$
  (GOE theory: 0.536)
\end{itemize}

\begin{remark}[Two Independent Results]
Section~\ref{sec:boundaries} answers ``why $1/2$?''
(CLT, universal). This section answers ``why GUE?''
($\CC$, specific). The first determines the \emph{location} of the
critical line; the second determines the \emph{fine structure}
of zeros on it.
\end{remark}

%=====================================================================
\section{Born-Rule Weighted Gram Graphs}
\label{sec:born}
%=====================================================================

\begin{definition}
For $N$ unit vectors $\psi_i \in \CC^d$, the \emph{Born-rule
weighted Gram graph} has adjacency
$W_{ij} = |\langle\psi_i|\psi_j\rangle|^2$ ($i \neq j$),
$W_{ii} = 0$.
\end{definition}

This uses the natural Born rule probability
$P = |\langle\psi|\phi\rangle|^2$
without artificial thresholding.

\subsection{Ramanujan condition}

A weighted graph is \emph{Ramanujan} if the second-largest
eigenvalue satisfies
$\lambda_2 \leq 2\sqrt{d_{\mathrm{eff}} - 1}$,
where $d_{\mathrm{eff}}$ is the average weighted degree.

\begin{theorem}[Near-Ramanujan, Numerical]\label{thm:ram}
For Born-rule weighted Gram graphs with $\rank \leq d$:
\begin{enumerate}[label=(\roman*)]
\item The Ramanujan ratio $\rho = \lambda_2 / 2\sqrt{d_{\mathrm{eff}}-1}$
  scales as $\rho \sim A \cdot d^{-\alpha}$ with
  $A \approx 1.94$, $\alpha \approx 0.67$ ($R^2 = 0.996$).
\item The critical dimension where $\rho = 1$ is
  $d_c \approx 3$.
\item For $d = 5$ (DRLT): $\rho < 1$ for all $N \leq N_c(5) \approx 500$.
  The graph is Ramanujan with $100\%$ probability
  up to $N \approx 300$.
\item For $N > N_c(d)$: $\rho$ exceeds $1$ by $O(1/\sqrt{N})$
  (soft boundary).
\end{enumerate}
\end{theorem}

\begin{center}
\begin{tabular}{lcccc}
\toprule
$d$ & $\rho$ ($N=100$) & Ramanujan\% & $N_c$ & $d > d_c$? \\
\midrule
3 & 0.96 & 77\% & $\sim 200$ & borderline \\
4 & 0.78 & 100\% & $\sim 300$ & Yes \\
\textbf{5} & \textbf{0.65} & \textbf{100\%} & $\sim$\textbf{500} & \textbf{Yes} \\
6 & 0.57 & 100\% & $\sim 700$ & Yes \\
10 & 0.41 & 100\% & $> 1000$ & Yes \\
\bottomrule
\end{tabular}
\end{center}

\begin{remark}
The critical dimension $d_c \approx 3$ is just below the
DRLT value $d = 5 = 2 + 3$ (sum of additive atoms).
This is not coincidental: the chiral decomposition
$\CC^5 = \CC^2 \oplus \CC^3$ guarantees $d > d_c$.
\end{remark}

\subsection{Ihara zeta zeros}

The Ihara zeta function of a graph encodes the ``primes'' of
the graph (primitive cycles).
For a Ramanujan graph, all Ihara zeros lie on the
critical line $|u| = 1/\sqrt{d-1}$, equivalent to
$\mathrm{Re}(s) = 1/2$.

\begin{proposition}[Discrete RH, Numerical]\label{prop:dRH}
For Born-rule weighted Gram graphs:
\begin{itemize}
\item $N \leq 50$, $d = 5$: $100\%$ of nontrivial Ihara zeros
  lie on $\mathrm{Re}(s) = 1/2$.
\item Primitive cycle counts follow the graph-PNT:
  $\pi(\ell) \sim (d-1)^\ell / \ell$.
\end{itemize}
\end{proposition}

%=====================================================================
\section{The Self-Contradiction Boundary}
\label{sec:boundary}
%=====================================================================

\begin{theorem}[Self-Contradiction]\label{thm:sc}
For the Gram ensemble $\mathcal{E}_N$ with $\Tr(G) = N$:
\begin{enumerate}[label=(\roman*)]
\item The resolution limit
  $\delta(N) = 1 - \max_{i \neq j} |G_{ij}|^2 > 0$
  for all finite $N$.
\item $\delta(N) = \Theta(N^{-1/2})$ ($R^2 = 0.9992$,
  EXP\_071).
\item $\delta(N) = 0$ requires $N = \infty$, which violates
  $\Tr(G) = N < \infty$.
\end{enumerate}
\end{theorem}

This creates a three-level picture:

\begin{center}
\begin{tabular}{lll}
\toprule
Level & Statement & Status \\
\midrule
$N \leq N_c(5)$ & Discrete RH: Ihara zeros on
  $\mathrm{Re}(s) = 1/2$ & \textbf{Theorem} \\
$N > N_c(5)$ & Zeros at $\mathrm{Re}(s) = 1/2 \pm O(1/\sqrt{N})$
  & \textbf{Theorem} \\
$N \to \infty$ & Exact $\mathrm{Re}(s) = 1/2$ (continuous RH)
  & Self-contradiction \\
\bottomrule
\end{tabular}
\end{center}

%=====================================================================
\section{Discussion}
\label{sec:disc}
%=====================================================================

\subsection{What this paper proves}

\begin{enumerate}
\item $\CC$ is the unique algebra where $\sigma_{\mathrm{stat}} =
  \sigma_{\mathrm{geom}} = 1/2$ (Theorem~\ref{thm:two}).
\item The number $2$ is uniquely doubly irreducible, explaining
  $n_T = \dim_\RR(\CC)$ (Theorem~\ref{thm:doubly}).
\item Born-rule Gram graphs are Ramanujan for $d \geq 4$
  up to $N_c(d)$ (Theorem~\ref{thm:ram}).
\item The discrete RH holds exactly for finite $N$
  (Proposition~\ref{prop:dRH}).
\end{enumerate}

\subsection{What this paper does not prove}

The Riemann Hypothesis itself remains open.
The gap between the discrete RH (finite Gram graphs)
and the continuous RH ($\zeta(s)$) is the self-contradiction
boundary: the continuum limit $N \to \infty$ is unreachable
within the DRLT axiom system ($\Tr(G) = N < \infty$).

The M\"obius randomness hypothesis --- that $\mu(n)$ behaves
like random unit-modulus coefficients --- would close this gap,
but is itself equivalent to RH.

\subsection{The (3,2) asymmetry}

The asymmetry between the spatial sector ($n_A = 3$) and
temporal sector ($n_T = 2$) is the asymmetry between:
\begin{itemize}
\item $3$: additive atom only (IVT: every cubic has a real root)
\item $2$: additive atom AND extension atom (doubly irreducible)
\end{itemize}
This asymmetry generates all of physics in DRLT,
and places $d = 5 = 2 + 3$ safely above the Ramanujan
critical point $d_c \approx 3$.

\subsection{Relation to existing programs}

\begin{center}
\begin{tabular}{lll}
\toprule
Program & What it explains & What we add \\
\midrule
Montgomery--Odlyzko & GUE $\leftrightarrow$ $\zeta$ zeros & Why GUE: $\CC \to \beta = 2$ \\
Hilbert--P\'olya & Self-adjoint operator & $G = \Psi\Psi^\dagger$ \\
Random matrix theory & Zero correlations & Why this ensemble \\
M\"obius randomness & $\mu(n)$ ``looks random'' & $\CC$-uniform phase \\
\bottomrule
\end{tabular}
\end{center}

%=====================================================================
% References
%=====================================================================
\begin{thebibliography}{99}

\bibitem{Paper1}
M.~Jeong and Claude,
``Chiral Atomic Decomposition of $\CC^5$,''
DRLT Paper Series (2026).

\bibitem{Paper2}
M.~Jeong and Claude,
``From Frobenius to Gauge: $\CC$ as the Unique Substrate,''
DRLT Paper Series (2026).

\bibitem{Paper3}
M.~Jeong and Claude,
``Zero-Parameter Predictions from the Gram Matrix,''
DRLT Paper Series (2026).

\bibitem{Paper4}
M.~Jeong and Claude,
``Coupling Running from Trace Decomposition:
$\zeta(s)$ as a Physical Parameter on $\CC^5$,''
DRLT Paper Series (2026).

\bibitem{Dyson}
F.~J.~Dyson,
``The threefold way: Algebraic structure of symmetry groups
and ensembles in quantum mechanics,''
J.\ Math.\ Phys.\ \textbf{3}, 1199 (1962).

\bibitem{Atas2013}
Y.~Y.~Atas, E.~Bogomolny, O.~Giraud, and G.~Roux,
``Distribution of the ratio of consecutive level spacings
in random matrix ensembles,''
Phys.\ Rev.\ Lett.\ \textbf{110}, 084101 (2013).

\bibitem{MontgomeryOdlyzko}
H.~L.~Montgomery,
``The pair correlation of zeros of the zeta function,''
Proc.\ Symp.\ Pure Math.\ \textbf{24}, 181 (1973);
A.~M.~Odlyzko, ``On the distribution of spacings between zeros
of the zeta function,''
Math.\ Comp.\ \textbf{48}, 273 (1987).

\bibitem{Harper}
A.~J.~Harper,
``Almost sure convergence of random multiplicative functions,''
arXiv:1301.2838 (2013).

\bibitem{Ihara}
Y.~Ihara,
``On discrete subgroups of the two by two projective linear
group over $\mathfrak{p}$-adic fields,''
J.\ Math.\ Soc.\ Japan \textbf{18}, 219 (1966).

\end{thebibliography}

\end{document}
